% !TeX root = ../教材题库_全部条目_转写工作稿.tex
% Source: XJ-B1-C05-EX-P160-Y0177
% 第5章 / 5.1 任意角和弧度制 / 教材 p.160 / PDF 第167页
% OCR未校对：公式与文字均可能有误；未完成数学TeX转写；禁止直接发给学生
% 原题及相邻原解截图（核对用）：../题目原图/XJ-B1-C05-EX-P160-Y0177-01.webp
% 本文件只在转写完成后作为题干使用；OCR 条目尚不是可靠题干。
\begin{bankitem}{5.1 例题 例题4}{第5章 / 5.1 任意角和弧度制 / 教材 p.160 / PDF 第167页}{XJ-B1-C05-EX-P160-Y0177}{OCR未校对：公式与文字均可能有误；未完成数学TeX转写；禁止直接发给学生}
\begin{ocrdraft}
[教材p.160 / PDF第167页]
\par
0互把下列各角从弧度化为度：
\par
rad；
\par
（2）5rad.
\par
4
\par
解（1）昼丕rad=—\ensuremath{\times}
\par
3丌180
\par
（2）5rad=5><
\par
286．5。．
\par
在科学应用中，用弧度制表示角的时候，单位“弧度”通常略去不写．比如，
\par
=1表示是1rad的角，即仅
\par
57。18'\ensuremath{\cdot}
\par
一2丌则表示=2丌rad，即是周角．因
\par
此，当直接用一个实数来表示角的时候，一定是指这个角的弧度数．
\par
在弧度制下，角的集合与实数集R之间建立起一一对应
\par
正角
\par
正实数
\par
的关系：每一个角都有唯一的一个实数（即这个角的弧度数）
\par
零角
\par
零
\par
与它对应；反过来，每一个实数也都有唯一的一个角（即弧
\par
负角
\par
负实数
\par
度数等于这个实数的角）与之对应，如图5．1一5所示．弧度
\par
角的集合
\par
实数集R
\par
数是实数，这将为我们今后用函数观点讨论涉及角的计算问
\par
图5．1一5
\par
题带来方便．
\par
利用弧度制度量角还有一个重要的原因，就是它能简化许多公式．例如，圆心
\par
角a一。时，所对弧长的计算公式是/=
\par
而根据弧度数的计算公式
\par
180
\par
若仅一工rad,得到弧长的另一计算公式：
\end{ocrdraft}
\end{bankitem}
